\documentclass[11pt]{article}

\def\chapitre{19}
\def\pagetitle{Équations différentielles.}

\input{setup.tex}

\begin{document}

\input{title.tex}

Une équation différentielle d'ordre $n\in\N^*$ est une équation qui exprime la dérivée $n^{\nt{ème}}$ d'une fonction en fonction des dérivées précédentes :
\begin{equation*}
    (E) ~ : ~ y^{(n)} = f(y,y', ..., y^{(n-1)})
\end{equation*}
avec $y:I\to F$ et $f:F^n \to F$ où $I$ est intervalle de $\R$ et $F$ un evdf.

Par définition, une solution de $E$ est une application $y:I\to F$, $n$ fois dérivables sur $I$ et telle que :
\begin{equation*}
    \forall t \in I, ~ y^{(n)}(t) = f(y(t), y'(t), ..., y^{(n-1)}(t)).
\end{equation*}
L'équation est dite linéaire si l'ensemble des solutions de $E$ est stable par combinaisons linéaires.

\section{Généralités sur les systèmes différentiels.}

\begin{defi}{}{}
    Soit $I$ un intervalle de $\R$, $F$ un $\R$-evdf, $a:I\to\L(F)$ continue sur $I$, $b:I\to F$ continue sur $I$.
    \begin{equation*}
        (E) ~ : ~ y ' = a(t)(y) + b(t)
    \end{equation*}
    Une solution de $E$ est une fonction $y:I\to F$ dérivable sur $I$ telle que $\forall t \in I, ~ y'(t)=a(t)(y(t))+b(t)$.\\
    \bf{Remarque.} $y$ est $\cursive{C}^1(I)$ comme combinaison linéaire.
\end{defi}

\begin{prop}{Traduction matricielle.}{}
    Soit $n$ la dimension de $F$, et $(e_1,...,e_n)$ une base de $F$. Pour tout $t\in I$, on pose :\\
    $\hra$ $B(t)$ le vecteur des coordonnées de $b(t)\in F$,\\
    $\hra$ $A(t)$ la matrice de $a(t)$ dans la base $(e_1,...,e_n)$,\\
    $\hra$ $X(t)$ le vecteur des coordonnées de $y(t)$.
    \begin{equation*}
        \forall t \in I, ~ y'(t) = a(t)(y(t)) + b(t) \iff \forall t \in I, ~ X'(t) = A(t)X(t) + B(t)
    \end{equation*}
    On notera $E_H$ l'équation différentielle homogène associée : $(E_H) ~ : ~ X'(t) = A(t)X(t)$.\\
    On notera $S_H$ l'ensemble des solutions de $E_H$, et $S$ l'ensemble des solutions de $E$.
\end{prop}

\begin{thm}{de Cauchy. (Admis)}{}
    Soit $I$ un intervalle de $\R$, $n\in\N^*$, $B:I\to\R^n$ et $A:I\to\M_n(\R)$ continues sur $I$.
    \begin{equation*}
        (E) ~ : ~ X' = A(t)X + B(t)
    \end{equation*}
    Pour $t_0\in I$ et $X_0\in\R^n$, il existe une unique solution $X$ de $E$ vérifiant $X(t_0)=X_0$.\\
    De plus, $X(t)=X_0+\int_{t_0}^t (A(s)X(s)+B(s))\nt{ds}$.
\end{thm}

\begin{prop}{Structure de $E_H$.}{}
    L'ensemble des solutions de $(E_H)$ a une structure d'espace vectoriel de \bf{dimension $n$}.\\
    Toute base de cet espace est appelé \bf{système fondamental de solutions}, qu'on notera (S.F.).
    \tcblower
    L'ensemble des solutions est un $\R$-ev (facile). Exhibons un isomorphisme entre l'espace des solutions et $\R^n$.\\
    Soit $t_0\in I$, $X\in\R^n$. D'après le théorème de Cauchy, il existe une unique solution de $(E_H)$ telle que $X(t_0)=X_0$.\\
    Posons alors $\phi:X\mapsto X(t_0)$, de $S_H$ dans $\R^n$, morphisme d'évaluation.\\
    Par le théorème de Cauchy, $\phi$ est bijective : c'est un isomorphisme donc $\dim S_H = \dim \R^n = n$.
\end{prop}

\begin{prop}{}{}
    Soit $I$ un intervalle de $\R$, $B:I\to\R^n$ et $A:I\to\M_n(\R)$ continues sur $I$.\\
    L'ensemble des solutions de $E$ est un espace affine de dimension $n$, de direction l'ensemble des solutions de $E_H$ :
    \begin{equation*}
        X \nt{ est solution de } (E) \iff \exists (X_P, X_H) \in S\times S_H, ~ X = X_P + X_H.
    \end{equation*}
    \tcblower
    Soit $X$ une solution de $E$ et $X_P$ une solution particulière (il en existe d'après le théorème de Cauchy).\\
    $\hra$ $X' = A(t)X + B(t)$ et $X'_P = A(t)X_P + B(t)$ pour tout $t \in I$.\\
    Alors $(X-X_P)' = A(t)(X-X_P) + \cancel{B(t) - B(t)}$ donc $X_H=X-X_P$ est solution de $E_H$.\\
    Réciproquement, si $X_P$ est solution de $E$ et $X_H$ solution de $E_H$, alors $X_P+X_H$ est solution de $E$.
\end{prop}

\begin{thm}{Principe de superposition.}{}
    Soient $B_1,B_2:I\to\R^n$ et $A:I\to\M_n(\R)$ continues sur $I$.
    \begin{align*}
        (E_1) ~ : ~ X'(t) = A(t)X(t) + B_1(t); \quad (E_2) ~ : ~ X'(t) = A(t)X(t) + B_2(t) 
    \end{align*}
    si $X_1$ est solution de $(E_1)$ et $X_2$ solution de $(E_2)$, alors $X_1+X_2$ est solution de $(E_1) + (E_2)$.
    \tcblower
    En effet, $(X_1+X_2)'=X_1' + X_2' = A(t)X_1+B_1(t) + A(t)X_2 + B_2(t) = A(t)(X_1+X_2) + B_1(t)+B_2(t)$.
\end{thm}

\begin{exercice}{}{}
    Soit $a\in\R$, $\phi:\R\to\R$ continue et périodique de période $T>0$.
    \begin{equation*}
        (E) ~:~ y' + \a y = \phi(t).
    \end{equation*}
    Soit $y$ une solution de $(E)$, montrer que $y$ est $T$-périodique ssi $y(0)=y(T)$.
    \tcblower
    D'après le théorème de Cauchy, $E$ admet des solutions.\\
    \boxed{\ra} Évident.\\
    \boxed{\la} On note $z(t)=y(t+T)$ pour $t\in\R$. Montrons que $z$ est solution de $(E)$.\\
    Déjà, $z$ est dérivable sur $\R$, puis $\forall t \in \R, ~ z'(t)+\a z(t) = y'(t+T)+\a y(t+T) = \phi(t+T) = \phi(t)$.\\
    De plus, $z(0)=y(T)=y(0)$, donc $y,z\in S$ et $y(0)=z(0)$, donc $y=z$ par théorème de Cauchy.
\end{exercice}

\vspace*{-0.8cm}

\begin{exercice}{}{}
    Pour $t\in\R$, on pose l'équation différentielle $(E)~:~(1+t^2)y'+2ty=1$.
    \begin{enumerate}
        \item Résoudre l'équation $(E)$.
        \item Résoudre le problème de Cauchy $y(1)=1$.
    \end{enumerate}
    \tcblower
    \boxed{1.} On a $\forall t \in \R, ~ 1+t^2\neq0$, on peut normaliser l'équation : $(E) : y' + \frac{2t}{1+t^2}y = \frac{1}{1+t^2}$.\\
    On pose $\forall t \in \R, ~ a(t) = \frac{2t}{1+t^2}$ et $b(t)=\frac{1}{1+t^2}$, continues sur $\R$.\\
    $\hra$ Solution de $E_H$ : $y_h=\l e^{-\ln|1+t^2|}=\frac{\l}{1+t^2}$ pour $\l\in\R$.\\
    $\hra$ Solution particulière : on pose $y(t)=\frac{\l(t)}{1+t^2}$, donc $y'(t)=\frac{(1+t^2)\l'(t)-2t\l(t)}{(1+t^2)^2}$ puis :
    \begin{equation*}
        \frac{(1+t^2)\l'(t)-2t\l(t)}{(1+t^2)^2} + \frac{2t\l(t)}{(1+t^2)^2} = \frac{1}{1+t^2}.
    \end{equation*}
    Il reste :
    \begin{equation*}
        \frac{\l'(t)}{1+t^2} = \frac{1}{1+t^2} \ra \l' = 1
    \end{equation*}
    Puis $\l(t)=t+\a$ avec $\a\in\R$, donc $y=\frac{1}{1+t^2} + \a\frac{1}{1+t^2}$ avec $\a\in\R$.\\
    \boxed{2.} $y(1)=1 \iff \frac{1}{2} + \frac{\a}{2} = 1 \iff \a = 1 \iff \forall t \in \R, ~ y(t) = \frac{1+t}{1+t^2}$.
\end{exercice}

\vspace*{-0.8cm}

\begin{exercice}{}{}
    Résoudre $y'=|y|$ sur $\R$.
    \tcblower
    Soit $y$ solution, alors $y'\geq0$ sur $\R$, donc $y$ est croissante sur $\R$.\\
    $\hra$ Supposons que $y$ ne s'annule pas, alors $y$ est continue sur $\R$ et ne s'annule pas : par TVI elle garde un signe constant puis :\\
    \hspace*{0.3cm} $\hra$ Si $y>0$,  $y=\l e^t$ pour $t\in\R$ et $\l>0$.\\
    \hspace*{0.3cm} $\hra$ Si $y<0$,  $y=\l e^{-t}$ pour $t\in\R$ et $\l<0$.\\
    $\hra$ Supposons que $y$ s'annule en $t_0\in\R$, elle est négative sur $]-\infty,t_0[$, positive sur $]t_0,+\infty[$.\\
    \hspace*{0.3cm} $\hra$ Sur $[t_0,+\infty[$ : $y=\l e^{t}$ avec $\l\geq 0$.\\
    \hspace*{0.3cm} $\hra$ Sur $]-\infty,t_0]$ : $y=\mu e^{-t}$ avec $\mu\leq 0$.\\
    Alors $y$ est bien définie sur $\R$ tout entier par morceaux, il faut vérifier que $y$ est dérivable sur $\R$.\\
    En particulier, $y$ est continue en $t_0$, $\l e^{t_0} = \mu e^{-t_0}=0$ donc $\l=\mu=0$ donc $y=0$.
\end{exercice}

\section{Exponentielles de matrices et d'endomorphismes.}

Rappelons que dans un espace vectoriel de dimension finie, une série absolument convergente est convergente.

\begin{defi}{}{}
    Soit $A\in\M_n(\K)$, on définit l'exponentielle de la matrice $A$ par : 
    \begin{equation*}\exp(A)=\sum_{k=0}^{+\infty}\frac{A^k}{k!}.\end{equation*}
    \tcblower
    Toutes les normes sont équivalentes sur $\M_n(\K)$, on choisit $\|\.\|$ une norme d'algèbre (existe) :\\
    $\hra$ $\forall A,B \in \M_n(\K), ~ \|AB\|\leq\|A\|\|B\|$ donc $\forall k \in \N, ~ \|A^k\|\leq\|A\|^k$.
    \begin{equation*}
        \forall k \in \N, ~ \left\|\frac{A^k}{k!}\right\| \leq \frac{\|A\|^k}{k!} ~ \nt{terme général d'une suite convergente réelle.}
    \end{equation*}
    Par théorème de comparaison des séries à terme général positif, $\sum\frac{A^k}{k!}$ converge absolument donc converge.
\end{defi}

\vspace*{-0.8cm}

\begin{ex}{}{}
    Pour une matrice nilpotente $A\in\M_n(\K)$ d'ordre $p\in\N$ :
    \begin{equation*}
        \exp(A) = \sum_{k=0}^{+\infty}\frac{A^k}{k!} = \sum_{k=0}^{p-1}\frac{A^k}{k!}
    \end{equation*} 
    Pour une matrice scalaire $A=\a I_n$ avec $\a\in\R$ :
    \begin{equation*}
        \exp(A) = \sum_{k=0}^{+\infty}\frac{(\a I_n)^k}{k!} = I_n\sum_{k=0}^{+\infty}\frac{\a^k}{k!} = e^\a I_n.
    \end{equation*}
    Pour une matrice diagonale $D=\Diag(\l_1,...,\l_n)$, $\exp(D)=\Diag(e^{\l_1},...,e^{\l_n})$.\\
    Pour une matrice diagonalisable $A=PDP^{-1}$, $\exp(A)=P\exp(D)P^{-1}$.
\end{ex}

\vspace*{-0.8cm}

\begin{prop}{}{}
    Soit $A\in\M_n(\K)$, alors $A$ et $\exp(A)$ commutent.
    \tcblower
    Rappel : $A$ commute avec tout polynôme en $A$ donc $\forall n \in \N, ~ A\left( \sum_{k=0}^n \frac{A^k}{k!}  \right) = \left( \sum_{k=0}^n\frac{A^k}{k!} \right)A$.\\
    Par continuité de $M\mapsto AM$ et $M\mapsto MA$ (linéaires en dim. finie) et par passage à la limite : $A\exp(A)=\exp(A)A$.
\end{prop}

\vspace*{-0.8cm}

\begin{prop}{}{}
    Soit $A\in\M_n(\K)$, alors $\exp(A)\in\K[A]$ : c'est un polynôme en $A$.
    \begin{equation*}
        \exists ! P \in \K_{d-1}[X], ~ \exp(A) = P(A) \nt{ où } d=\deg\Pi_A
    \end{equation*}
    \tcblower
    Puisque $\K[A]$ est fermé comme $\K$-algèbre en dimension finie, par passage à la limite : $\exp(A)\in\K[A]$.\\
    Soit $P\in\K_{d-1}[X]$ tel que $\exp(A)=P(A)$, on a $\exp(A)=\sum_{k=0}^{+\infty}\frac{A^k}{k!} = P(A)$.\\
    \bf{Unicité.} Pour $P(A)=Q(A)$ avec $\deg(P)<d$ et $\deg(Q)<d$, on a $(P-Q)(A)=0$ donc $P-Q=0$ par déf de $d$.
\end{prop}

\begin{ex}{}{}
    Soit $D=\Diag(\l_1,...,\l_p)$. Calculer le polynôme $Q$ tel que $\exp(D)=Q(D)$.
    \tcblower
    Soient $L_1,...,L_p$ les polynômes de Lagrange associés : $\forall i,j : L_i(\l_j)=\d_{i,j}$.\\
    Ces polynômes forment une base de $\R_{p-1}[X]$ puis $\forall Q \in \R_{p-1}[X], ~ Q = \sum_{i=1}^pQ(\l_i)L_i$.\\
    Soit $Q$ l'unique polynôme de $\R_{p-1}[X]$ tel que $\forall i \in \lb1,p\rb, ~ Q(\l_i)=e^{\l_i}$ : $Q=\sum_{i=1}^pe^{\l_i}L_i$.\\
    Alors $Q(D)=\exp(D)$, avec $Q$ qui est bien de degré inférieur à $\Pi_D$.
\end{ex}

\vspace*{-0.8cm}

\begin{prop}{}{}
    Soient $A,B\in\M_n(\K)$ tels que $A=PBP^{-1}$ avec $P\in\GL_n(\K)$, alors $\exp(A)=P(\exp B)P^{-1}$.
    \tcblower
    Par récurrence sur $k\in\N$ : $\forall k \in \N, ~ A^k = PB^kP^{-1}$ puis :
    \begin{equation*}
        \forall n_0 \in \N, ~ \sum_{k=0}^{n_0} \frac{A^k}{k!} = \sum_{k=0}^{n_0}P\frac{B^k}{k!}P^{-1} = P\left( \sum_{k=0}^{n_0}\frac{B^k}{k!} \right)P^{-1}
    \end{equation*}
    Puis par continuité de $M\mapsto PMP^{-1}$ (linéaire en dim. finie) et par passage à la limite : $\exp(A)=P\exp(B)P^{-1}$.
\end{prop}

\vspace*{-0.8cm}

\begin{exercice}{}{}
    Soit :
    \begin{equation*}
        A=\begin{pmatrix}0&1&1\\1&0&1\\1&1&0\end{pmatrix} \in \S_3(\R).
    \end{equation*}
    \begin{enumerate}
        \item Déterminer $\X_A$ et $\Pi_A$.
        \item Montrer que $A$ est diagonalisable, la diagonaliser et en déduire $\exp(A)$.
        \item Calculer $A^n$ pour tout $n\in\N$.
        \item Retrouver $\exp(A)$.
    \end{enumerate}
    \tcblower
    \boxed{1.} $A$ est diagonalisable par théorème spectral. Soit $\l\in\R$. $C_1\gets C_1 + C_2 + C_3$; $L_2 \gets L_2 - L_1$ et $L_3\gets L_3 - L_1$.
    \begin{equation*}
        \det(\l I_n - A) = \begin{vmatrix} \l & -1 & -1 \\ -1 & \l & -1 \\ -1 & -1 & \l\end{vmatrix} = (\l-2)\begin{vmatrix}1 & -1 & -1 \\ 1 & \l & -1 \\ 1 & -1 & \l\end{vmatrix} = (\l-2)\begin{vmatrix}1&-1&-1\\0&\l+1&0\\0&0&\l+1\end{vmatrix}=(\l-2)(\l+1)^2.
    \end{equation*}
    Donc $\X_A = (X-2)(X+1)^2$. Puis $\Pi_A=(X+1)(X+2)$ (car SRS et les vp sont exactement racines de $\Pi_A$).\\
    \boxed{2.} $E_2=\Vect(1,1,1)$, $E_{-1}=\Vect((\frac{1}{\sqrt{2}}, -\frac{1}{\sqrt{2}}, 0), (\frac{1}{\sqrt{6}}, \frac{1}{\sqrt{6}}, -\frac{\sqrt{2}}{\sqrt{3}}))$ donc $A=PDP^\top$ avec \begin{equation*}P=\begin{pmatrix}\frac{1}{\sqrt{3}} & \frac{1}{\sqrt{2}} & \frac{1}{\sqrt{6}} \\ \frac{1}{\sqrt{3}} & -\frac{1}{\sqrt{2}} & \frac{1}{\sqrt{6}} \\ \frac{1}{\sqrt{3}} & 0 & -\frac{\sqrt{2}}{\sqrt{3}}\end{pmatrix} \et D=\begin{pmatrix}2&0&0\\0&-1&0\\0&0&-1\end{pmatrix} \ra \exp(A) = P(\exp D)P^\top = \frac{1}{3e}\begin{pmatrix}  2+e^3 & e^3-1 & e^3-1 \\ e^3-1 & 2+e^3 & e^3-1 \\ e^3-1 & e^3-1 & 2+e^3 \end{pmatrix}\end{equation*}
    \boxed{3.} $X^n = \Pi_A(X)Q(X)+R_n(X)$ avec $R_n(X)=a_n X + b_n I_3$ par division euclidienne donc $A^n=R_n(A)$.\\
    En $X=2$ : $2^n = 2 a_n + b_n$; en $X=-1$ : $(-1)^n = -a_n+b_n$ donc $a_n = \frac{2^n}{3}+\frac{(-1)^{n+1}}{3}$ et $b_n = \frac{2^n}{3}+\frac{2(-1)^n}{3}$.
\end{exercice}

\begin{prop}{}{}
    Soient $A,B\in\M_n(\K)$ telles que $AB=BA$, alors $\exp(A)\exp(B)=\exp(A+B)$.
    \tcblower
    Les séries $\exp(A)$ et $\exp(B)$ sont absolument convergentes, donc par produit de Cauchy :
    \begin{equation*}
        \left( \sum_{k=0}^{+\infty}\frac{A^k}{k!} \right)\left( \sum_{k=0}^{+\infty}\frac{B^k}{k!} \right) = \sum_{n=0}^{+\infty}\sum_{k=0}^n \frac{A^k}{k!}\frac{B^{n-k}}{(n-k)!} = \sum_{n=0}^{+\infty} \frac{1}{n!}\sum_{k=0}^{+\infty}\binom{n}{k}A^kB^{n-k} = \sum_{n=0}^{+\infty}\frac{(A+B)^n}{n!}
    \end{equation*}
    Donc $\exp(A)\exp(B)=\exp(A+B)$. La commutativité est utilisée par le binôme de Newton.
\end{prop}

\vspace*{-0.9cm}

\begin{prop}{}{}
    Soit $A\in\M_n(\K)$, alors $\exp(A)\in\GL_n(\K)$ et $(\exp A)^{-1} = \exp(-A)$.
    \tcblower
    \small $A$ et $-A$ commutent donc $\exp(A)\exp(-A)=\exp(A-A)=\exp(0)=I_n$ donc $\exp(A)\in\GL_n(\K)$ et $(\exp A)^{-1}=\exp(-A)$.
\end{prop}

\vspace*{-0.9cm}

\begin{prop}{}{}
    Soit $A\in\M_n(\K)$ telle que $\X_A$ est scindé, alors $\det(\exp(A))=e^{\Tr(A)}$.
    \tcblower
    Puisque $\X_A$ est scindé, $A$ est trigonalisable : $\exists P \in \GL_n(\K), ~ \exists T \in \T_n^+(\K), ~ A=PTP^{-1}$.\\
    On rappelle que la diagonale de $T$ est composée des valeurs propres $(\l_1,...,\l_n)$ de $A$, et $\exp(A)=P\exp(T)P^{-1}$.
    \begin{equation*}
        \det(\exp A) = \det(\underbrace{\exp T}_{\in\T_n^+(\K)}) = \prod_{i=1}^n \exp(\l_i) = \exp(\sum_{i=1}^n \l_i) = e^{\Tr(A)}
    \end{equation*}
    En effet, puisque $\X_A$ est scindé, on a $\sum_{i=1}^n \l_i = \Tr(A)$, et on retrouve que $\exp(A)\in\GL_n(\K)$.
\end{prop}

\vspace*{-0.9cm}

\begin{exercice}{}{}
    Montrer que $A\in\A_n(\R) \ra \exp(A) \in O_n(\R)$, et même $A\in\SO_n(\R)$.\\
    Pour $A\in\A_2(\K)$, donner la forme de $\exp(A)$.
    \tcblower
    \small
    Soit $A\in\A_n(\R)$ : $A^\top = -A$.\\
    $\hra$ Déjà, $\forall M \in \M_n(\R), ~ \exp(M)^\top = \exp(M^\top)$ :
    \begin{equation*}
        \forall n_0 \in \N, ~ \left( \sum_{k=0}^{n_0}\frac{M^k}{k!} \right)^\top = \sum_{k=0}^{n_0}\frac{(M^\top)^k}{k!}
    \end{equation*}
    On a utilisé la lin. de $M\mapsto M^\top$; puis par continuité de $M\mapsto M^\top$, lin. en dim. finie : $\exp(M)^\top = \exp(M^\top)$.\\
    On revient à $A$ : $(\exp A)^\top = \exp(A^\top) = \exp(-A)=(\exp A)^{-1}$, donc $\exp A \in O_n(\R)$.\\
    De plus, $\det(\exp A) = e^{\Tr(A)} > 0$ donc $\exp(A)\in\SO_n(\R)$, en fait $\det(\exp A)=1$.\\
    $\hra$ Soit $A\in\A_2(\R)$, $\exists \a \in \R$ tel que :
    \begin{equation*}
        A = \begin{pmatrix}0 & \a \\ -\a & 0 \end{pmatrix}; ~ \forall n \in \N, ~ A^{2n} = (-1)^n\a^{2n}I_2; ~ A^{2n+1} = (-1)^n\a^{2n}A
    \end{equation*}
    Par sommation par paquets :
    \begin{equation*}
        \exp A = \sum_{n=0}^{\infty}\frac{A^n}{n!} = \sum_{n=0}^{\infty}\frac{A^{2n}}{(2n)!} + \sum_{n=0}^{+\infty}\frac{A^{2n+1}}{(2n+1)!} = \sum_{n=0}^\infty\frac{(-1)\a^{2n}}{(2n)!}I_2 + \sum_{n=0}^{+\infty}\frac{(-1)^n\a^{2n}}{(2n+1)!}A = \cos(\a)I_2 + \frac{1}{\a}\sin(\a)A
    \end{equation*}
\end{exercice}

\begin{exercice}{}{}
    Soit $A\in\M_n(\C)$. Montrer que $\ds\lim_{p\to+\infty}\left( I_n + \frac{A}{p} \right)^p = \exp(A)$.
\end{exercice}

\begin{prop}{}{}
    Soit $A\in\M_n(\K)$, on pose $\phi:t\mapsto\exp(tA)$, alors $\phi$ est dérivable sur $\R$ et $\forall t \in \R, ~ \phi'(t)=A\exp(tA)=\exp(tA)A$.
    \tcblower
    Soit $t_0\in\R$, $\forall h \in \R, ~ \exp((t_0 + h)A) - \exp(t_0A) = \exp(t_0A)(\exp(hA) - I_n)$.\\
    Posons $\Delta_h = \frac{\phi(t_0 + h) - \phi(t_0)}{h}$ pour tout $h\neq0$.
    \begin{equation*}
        \Delta_h - \exp(t_0A) = \frac{\exp((t_0 + h)A) - \exp(t_0A)}{h} - \exp(t_0A)A = \exp(t_0A)\left(\frac{\exp(hA) - I_n}{h} - 1\right)
    \end{equation*}
    Or : 
    \begin{equation*} \frac{\exp(hA) - I_n}{h} - A = \frac{1}{h}\left( \sum_{k=0}^\infty \frac{h^kA^k}{k!} - I_n \right)-A = \frac{1}{h}\left( \sum_{k=1}^{\infty}\frac{h^kA^k}{k!} \right) - A = \sum_{k=2}^{\infty}\frac{h^{k-1}A^k}{k!}\end{equation*}
    On choisit une norme sous-multiplicative $\|\.\|$ (toutes les normes sont équivalentes).
    \begin{align*}
        \left\|\frac{\exp(hA) - I_n}{h}\right\| &= \left\| \sum_{k=2}^{\infty}\frac{h^{k-1}A^k}{k!} \right\| \leq \sum_{k=2}^{\infty}\frac{|h|^{k-1}\|A\|^k}{k!} = \frac{1}{|h|}\sum_{k=2}^{\infty}\frac{(|h|\|A\|)^k}{k!} \\&= \frac{1}{|h|}\left( e^{|h|\|A\|} - |h|\|A\| - 1 \right) \sim_0 \frac{1}{|h|}\frac{(|h|\|A\|)^2}{2} \xrightarrow[h\to0]{}0.
    \end{align*}
    Par encadrement, $\left\|\frac{\exp(hA)-I_n}{h}\right\|\xrightarrow[h\to0]{}0$ puis $\|\Delta_h-\exp(t_0A)A\|\xrightarrow[h\to0]{}0$ i.e. $\Delta_h \xrightarrow[h\to0]{}\exp(t_0A)A$.\\
    Donc $\phi$ est dérivable en $t_0$ et $\phi'(t_0)=\exp(t_0A)A$ puis $\exp(t_0A)$ est un polynôme en $A$ donc commute avec $A$.
\end{prop}

\section{Systèmes différentiels linéaires autonomes.}

Un système différentiel linéaire d'ordre 1 est dit autonome si $A$ ne dépend pas de $t$.

\begin{thm}{Résolution de \texorpdfstring{$X'=AX$}{Lg}}{}
    Soit $A\in\M_n(\K)$ et $(E) ~ : ~ X' = AX$.\\
    Alors $X$ est solution de $(E)$ si et seulement si $\exists C \in \M_{n,1}(\K), ~ \forall t \in \R, ~ X(t) = \exp(tA)C$.
    \tcblower
    Soit $X:\R\to\M_{n,1}(\K)$ et posons $Y(t)=\exp(-tA)X(t)$ pour tout $t\in\R$. $X$ est dérivable ssi $Y$ l'est et :
    \begin{equation*}
        \forall t \in \R, ~ Y'(t) = -A\exp(tA)X(t)+\exp(-tA)X'(t)
    \end{equation*}
    Donc, puisque $X'=AX$ :
    \begin{equation*}
        X' = AX \iff Y' = -A\exp(-tA)X + \exp(-tA)X' = \exp(-tA)AX - \exp(-tA)AX = 0 \iff Y = \cste
    \end{equation*}
    Donc $X'=AX\iff X=(\exp tA)Y = \exp(tA)C$ avec $C\in\M_{n,1}(\K)$.
\end{thm}

\begin{prop}{}{}
    Soit $A\in\M_n(\K)$ et $(E) ~:~ X'=AX$. L'ensemble $S$ des solutions de $(E)$ est un $\K$-ev de dimension $n$.\\
    Un système fondamental de solutions est donné par les vecteurs colonnes de la matrice $\exp(tA)$.
    \tcblower
    On a déjà vu que l'espace vectoriel $S$ des solutions de $E$ est un $\K$-ev de dimension $n$.
    \begin{equation*}
        X' = AX \iff \exists C \in \M_{n,1}(\K), ~ \forall t \in \R, ~ X(t)=(\exp tA)C.
    \end{equation*}
    Notons $X_1(t),...,X_n(t)$ les vecteurs colonnes de $\exp(tA)$ et $C=(\a_i)\in\M_{n,1}(\K)$, alors :
    \begin{equation*}
        X(t) = \begin{pmatrix}x_{1,1}(t) & \dots & x_{1,n}(t) \\ \vdots & & \vdots \\ x_{n,1}(t) & \dots & x_{n,n}(t) \end{pmatrix}\begin{pmatrix}\a_1 \\ \vdots \\ \a_n\end{pmatrix} = \begin{pmatrix}\a_1 x_{1,1}(t) + \dots + \a_n x_{1,n}(t) \\ \vdots \\ \a_1x_{n,1}(t) + \dots + \a_nx_{n,n}(t) \end{pmatrix}= \a_1X_1(t) + ... + \a_nX_n(t)
    \end{equation*}
    Finalement, $X=\a_1X_1+...+\a_nX_n$, donc $(X_1,...,X_n)$ est génératrice de $S$, de cardinal $n$, donc une base de $S$.
\end{prop}

\vspace*{-0.9cm}

\begin{thm}{}{}
    Soit $A\in\M_n(\K)$ \bf{diagonalisable}, avec $\Sp(A)=\{\l_1,...,\l_n\}$ éventuellement répétées.
    On note $(V_1,...,V_n)$ une base de vecteurs propres de $A$, et $(E) ~:~ X'=AX$.\\
    Un système fondamental de solution est donné par $(X_1,...,X_n)$ où $\forall i \in \lb1,n\rb, ~ \forall t \in \R, ~  X_i(t)=e^{\l_i t}V_i$.
    \tcblower
    Diagonalisons $A$ avec $A=PDP^{-1}$ : $D=\Diag(\l_1,...,\l_n)$ et $P=(V_1 \dots V_n)$.\\
    On a $X'=AX\iff X'=PDP^{-1}X$, on pose $Y=P^{-1}X$, $Y$ est dérivable sur $\R$ et $Y'=P^{-1}X'$.\\
    Donc $X'=AX \iff P^{-1}X'=D(P^{-1}X)\iff Y'=DY \iff \forall i \in \lb1,n\rb, ~ y_i' = \l_iy_i \iff y_i = \a_i e^{\l_it}$.\\
    Alors $\forall t \in \R, ~ Y(t)=(\a_i e^{\l_i t})\in\M_{n,1}(\K)$ avec $\a_1,...,\a_n \in \K$.\\
    Enfin, $X=PY$ donc $X=\a_1e^{\l_1 t}V_1 + ... + \a_1e^{\l_nt}V_n$. On pose $\forall i \in \lb1,n\rb, ~ X_i = e^{\l_i t}V_i$.\\
    Donc $(X_1,...,X_n)$ est génératrice de $S$ de cardinal $n$, c'est une base de $S$.
\end{thm}

\vspace*{-0.9cm}

\begin{exercice}{}{}
    Résoudre :
    \begin{equation*}
        \begin{cases}
            x' = 2x - 2y \\ y' = -x + 3y
        \end{cases}
    \end{equation*}
    \tcblower
    Écriture matricielle : $X'=AX$ avec $X=\begin{pmatrix}x\\y\end{pmatrix}$ et $A=\begin{pmatrix}2&-2\\-1&3\end{pmatrix}$.\\
    L'ensemble des solutions est un $\R$-ev de dimension 2. On a $\X_A = (X-1)(X-4)$, SRS donc $A$ est diagonalisable.\\
    Les valeurs propres de $A$ sont $1$ et $4$ simples et $E_1=\Vect(2,1)$ et $E_4=\Vect(1,-1)$.\\
    Par théorème, $(X_1,X_2)$ est un système fondamental de solutions avec :
    \begin{equation*}
        X_1(t) = e^{\l_1t}V_1 = e^t\begin{pmatrix}2\\1\end{pmatrix} = \begin{pmatrix}2e^t\\e^t\end{pmatrix} \et X_2(t) = e^{\l_2 t}V_2 = \begin{pmatrix}e^{4t}\\-e^{4t}\end{pmatrix}
    \end{equation*}
    Puis $X$ solution de $X'=AX$ ssi $\exists (\a_1,\a_2)\in\R^2, ~ X=\a_1X_1+\a_2X_2$ ssi $\forall t \in \R, ~ X(t)=\begin{pmatrix}2\a_1e^t + \a_2e^{4t} \\ \a_1e^t - \a_2e^{4t}\end{pmatrix}$.\\
    Conclusion : $x$ et $y$ sont solutions ssi $\exists (\a_1,\a_2)\in\R^2$ tels que :
    \begin{equation*}
        \begin{cases}
            x(t) = 2\a_1e^t + \a_2e^{4t}\\
            y(t) = \a_1e^t - \a_2e^{4t}
        \end{cases}
    \end{equation*}
\end{exercice}

\vspace*{-0.8cm}

\begin{exercice}{}{}
    Résoudre :
    \begin{equation*}
        \begin{cases}
            x' = x - y \\ y' = x + y
        \end{cases}
    \end{equation*}
    Avec $x(0)=0$ et $y(0)=1$.
    \tcblower
    Écriture matricielle : $X'=AX$ avec $X=\begin{pmatrix}x\\y\end{pmatrix}$ et $A=\begin{pmatrix}1&-1\\1&1\end{pmatrix}$.\\
    On reconnaît un système différentiel linéaire autonome homogène, on sait que l'ensemble de ses solution est un $\R$-ev de dimension 2 car $A\in\M_2(\R)$. Le théorème de Cauchy s'applique : il existe une unique solution vérifiant la condition initiale donnée.\\
    $\hra$ On regarde si $A$ est diagonalisable : $\X_A=X^2-2X+2=(X-1-i)(X-1+i)$, SRS dans $\M_2(\C)$, donc $A$ est dz dans $\M_2(\C)$. On a $E_{1+i}=\Vect(i, 1)$ et $E_{1-i}=\Vect(-i,1)$. On pose $V_1=(i~1)$ et $V_2=(-i~1)$.\\
    Par théorème, $X_1=e^{(1+i)t}V_1$ et $X_2=e^{(1-i)t}V_2$ forment un système fondamental de solutions.\\
    Alors $\exists!(\l,\mu)\in\C^2, ~ X = \l X_1 + \mu X_2 ~ : ~ X(t) = \l e^{(1+i)t}\begin{pmatrix}i\\1\end{pmatrix}+\mu e^{(1-i)t}\begin{pmatrix}-i\\1\end{pmatrix}$.\\
    Puis $X(0)=\begin{pmatrix}0\\1\end{pmatrix}=\l\begin{pmatrix}i\\1\end{pmatrix}+\mu\begin{pmatrix}-i\\1\end{pmatrix}$ donc $\l=\mu=\frac{1}{2}$ puis $X(t)=\frac{1}{2}e^{(1+i)t}\begin{pmatrix}i\\1\end{pmatrix}+\frac{1}{2}e^{(1-i)t}\begin{pmatrix}-i\\1\end{pmatrix}$.\\
    On sait cependant que l'unique solution est réelle, donc on peut travailler l'expression :
    \begin{equation*}
        \begin{cases}
            x(t) = \frac{1}{2}e^{(1+i)t}i - \frac{1}{2}e^{(1-i)t}i = e^t\frac{i}{2}(e^{it}-e^{-it})\\
            y(t) = \frac{1}{2}e^{(1+i)t} + \frac{1}{2}e^{(1-i)t} = \frac{e^t}{2}(e^{it} + e^{-it})
        \end{cases} \iff \begin{cases}x(t) = e^t \frac{i}2{(2i\sin t)} = -e^t\sin t \\ y(t) = \frac{e^t}{2}(2\cos t) = e^t \cos t\end{cases}
    \end{equation*}
\end{exercice}

\vspace*{-0.8cm}

\begin{exercice}{}{}
    Résoudre :
    \begin{equation*}
        \begin{cases}
            x' = 2x + z \\ y' = 2y - z \\ z' = x - y + z
        \end{cases}
    \end{equation*}
    \tcblower
    Traduction matricielle : on pose $X=\begin{pmatrix}x\\y\\z\end{pmatrix}$ et $A=\begin{pmatrix}2 & 0 & 1 \\ 0 & 2 & -1 \\ 1 & -1 & 1\end{pmatrix}$.\\
    Puisque $A$ est symétrique réelle, elle est orthodiagonalisable. $C_1\gets C_1 + C_2$ puis selon colonne 1.\\
    On a $\X_A(\l) = \l(\l-2)(\l-3)$, $E_0 = \Vect(1,-1,-2)$, $E_2 = \Vect(1,1,0)$, $E_3=\Vect(1,-1,1)$.
    \begin{equation*}
        A=PDP^{-1}; ~ D=\begin{pmatrix}0&0&0\\0&2&0\\0&0&3\end{pmatrix} \et P=\begin{pmatrix}1&1&1\\-1&1&-1\\-2&0&1\end{pmatrix}\in\GL_3(\R).
    \end{equation*}
    Un système fondamental de solutions est donné par $\begin{cases}X_1(t)=e^{0t}V_0 = V_0 \\ X_2(t)=e^{2t}V_1 \\ X_3(t)=e^{3t}V_2\end{cases}$.\\
    Puis $\exists!(\l,\mu,\g)\in\R^3, ~ X(t)=\l X_1(t)+\mu X_2(t) + \g X_3(t)$ i.e. 
    \begin{equation*}
        \begin{cases}
            x(t) = \l + \mu e^{2t} + \g e^{3t} \\ y(t) = -\l + \mu e^{2t} - \g e^{3t} \\ z(t) = -2\l + \g e^{3t}.
        \end{cases}
    \end{equation*}
\end{exercice}

Soit $X'(t)=A(t)X(t)+B(t)$ pour $t\in I$ un intervalle de $\R$ tels que $A:I\to\M_n(\K)$ et $B:I\to\M_{n,1}(\K)$ continues sur $I$.\\
Soit $(E_H) : ~ X'(t)=A(t)X(t)$ système différentiel homogène associé.

\begin{defi}{}{}
    Si $(X_1,...,X_n)$ est une famille de solutions de $(E_H)$, on pose $\forall t \in I, ~ W(t)=\det(X_1(t),...,X_n(t))$.\\
    Alors $W$ est le Wronskien de la famille $(X_1,...,X_n)$.
\end{defi}

\begin{prop}{}{}
    Soit $(X_1,...,X_n)$ une famille de solutions de $(E_H)$, les propositions suivantes sont équivalents :
    \begin{enumerate}
        \item $(X_1,...,X_n)$ est un S.F. de $E$.
        \item $\forall t \in I, ~ W(t) \neq 0$.
        \item $\exists t \in I, ~ W(t) \neq 0$.
    \end{enumerate}
    \tcblower
    \boxed{1\ra2} Soit $t_0\in I$ et $\phi:X\mapsto X(t_0)$ de $S$ dans $\K^n$.\\
    Par théorème de Cauchy, $\phi$ est bijection donc l'image par $\phi$ d'une base $(X_1,...,X_n)$ de $S$ est base de $\K^n$.\\
    Or $\phi(X_1,...,X_n)=(X_1(t_0),...,X_n(t_0))$ puis $(X_1(t_0),...,X_n(t_0))$ est une base de $\K^n$ ssi $\det(X_1(t_0),...,X_n(t_0))\neq0$ i.e. $W(t_0)\neq0$.\\
    \boxed{2\ra3} Immédiat.\\
    \boxed{3\ra 1} Par hypothèse, $\exists t_0 \in I, ~ W(t_0)\neq0$, par hypothèse, $(X_1(t_0),...,X_n(t_0))$ est base de $\K^n$.\\
    L'image par $\phi^{-1}$ de cette base est donc une base de $S$. Par définition, $(X_1,...,X_n)$ est donc un $S.F.$ de $E$.
\end{prop}

Pour trouver une solution particulière, on peut utiliser la méthode de variation des constantes :

\begin{prop}{}{}
    Soit $(E) ~:~ X'(t)=A(t)X(t)+B(t)$ avec $A:I\to\M_n(\K)$ et $B:I\to\M_{n,1}(\K)$ continues sur $I$.\\
    Soit $(X_1,...,X_n)$ un (S.F.) de solutions de $E$, i.e. de $E_H$.
    Posons $X(t) = \l_1(t)X_1(t)+...+\l_n(t)X_n(t)$ pour $t\in I$ avec $\l_i : I \to \K$ pour tout $i\in\lb1,n\rb$.\\
    Alors $X$ est dérivable sur $I$ ssi $\l_i$ est dérivable sur $I$ puis $X$ solution de $E$ ssi $\sum_{i=1}^n\l_i'(t)X_i(t)=B_i(t)$.\\
    Ceci est un système de Cramer qui donne de manière unique $\l'_1,...,\l_n'$.\\
    On les calcule avec $n$ calculs de primitives, alors $\l_i(t) = \a_i(t) + K_i$, avec $K_i$ une constante, puis :
    \begin{equation*}
        X(t) = \sum_{i=1}^n K_i X_i(t) + \sum_{i=1}^n \a_i(t)X_i(t).
    \end{equation*}
    \tcblower
    Cas particulier $n=1$ :\\
    $\hra$ $(E) ~:~ y'(t)=a(t)y(t)+b(t)$, $a,b$ continues de $I$ dans $\K$.\\
    $\hra$ $(E_H) ~:~ y'(t)=a(t)y(t)$.\\
    Les solutions de $E_H$ sont de la forme $y=\l e^{\a(t)}$ avec $\a$ primitive de $a$.\\
    Alors $y_0(t) = e^{\a(t)}$ où $y_0$ est un S.F de $E_H$ i.e. de $E$. Posons $y=\l(t)e^{\a(t)}$.\\
    On a $y$ dérivable sur $I$ ssi $\l$ est dérivable sur $I$ puis :
    \begin{align*}
        y \nt{ solution de } E &\iff \l'(t)e^{\a(t)} + \cancel{\l(t)\a(t)e^{\a(t)}} = \cancel{a(t)(\l(t)e^{\a(t)}) + b(t)}\\
        &\iff \l'(t)y_0(t) = b(t) \iff \l'(t)=b(t)e^{-\a(t)} \iff \l(t) = \mu(t) + K \\
        &\iff y(t) = (\mu(t) + K)e^{\a(t)} = \mu(t) e^{\a(t)} + Ke^{\a(t)}.
    \end{align*}
\end{prop}

Dans le cas général, en reprenant les notations de l'énoncé : $X(t)=\sum_{i=1}^n \l_i(t)X_i(t)$ avec $X_1,...,X_n$ S.F. de $E$.
Vérifions $X$ dérivable sur $I$ ssi $\forall i \in \lb1,n\rb, ~ \l_i$ dérivable sur $I$ :\\
\boxed{\la} Par opérations.\\
\boxed{\ra} Par hypothèse, $\l_1(t)X_1(t) + ... + \l_n(t)X_n(t)=X(t)$, i.e. 
\begin{equation*}
    \l_1(t) \begin{pmatrix}x_{1,1}(t) \\ \vdots \\ x_{n,1}(t)\end{pmatrix} + ... + \l_n(t)\begin{pmatrix}x_{1,n}(t) \\ \vdots \\ x_{n,n}(t)\end{pmatrix} = \begin{pmatrix}x_1(t) \\ \vdots \\ x_n(t)\end{pmatrix}
\end{equation*}
i.e.
\begin{equation*}
    \underbrace{\begin{pmatrix}
        x_{1,1}(t) & \dots & x_{1,n} \\ \vdots & & \vdots \\ x_{n,1}(t) & \dots & x_{n,n}(t)
    \end{pmatrix}}_{M(t)} \begin{pmatrix} \l_1(t) \\ \vdots \\ \l_n(t) \end{pmatrix} = \begin{pmatrix} x_1(t) \\ \vdots \\ x_n(t) \end{pmatrix}
\end{equation*}
On a $\det M(t) = W(t)$ (Wronskien) puis $\forall t \in I, ~ W(t)\neq 0$ par théorème donc $M(t)\in\GL_n(\K)$. Donc :
\begin{equation*}
    \begin{pmatrix}\l_1(t) \\ \vdots \\ \l_n(t) \end{pmatrix} = M(t)^{-1} \begin{pmatrix}x_1(t) \\ \vdots \\ x_n(t)\end{pmatrix}
\end{equation*}
Or $A^{-1}=\frac{1}{\det(A)}\Com(A)^\top$ si $A\in\GL_n(\K)$, donc $\forall t \in I, ~ (M(t))^{-1}=\frac{1}{W(t)}(\Com M(t))^\top$.\\
Puis $W$ est dérivable avec $W'(t)=\det(X_1'(t),X_2,...,X_n(t))+...+\det(X_1(t),...X_{n-1}(t),X_n'(t))$ et ne s'annulle pas.\\
Puis $t\mapsto M$ est dérivable sur $I$, puis $t\mapsto \Com M(t)$ est dérivable car $[\Com A]_{i,j}=(-1)^{i+j}\Delta_{i,j}$.\\
Puis $t\mapsto (\Com M)^\top$ est dérivable sur $I$, finalement $t\mapsto M(t)^{-1}$ est dérivable sur $I$ et $\forall i \in \lb1,n\rb$ dérivable sur $I$.\\
Enfin, soit $X$ d. sur $I$ tq $X=\sum_{i=1}^n\l_iX_i$, alors $\l_1,...,\l_n$ sont d. sur $I$ et $X'=\sum_{i=1}^n\l_i'X_i + \sum_{i=1}^n\l_i X_i'$ (Leibniz)
\begin{align*}
    X' = AX + B &\iff \sum_{i=1}^n \l_i'X_i + \sum_{i=1}^n\l_iX'_i = A\left( \sum_{i=1}^n \l_i X_i \right) + B \\
    &\iff \sum_{i=1}^n \l_i' X_i + \cancel{\sum_{i=1}^n \l_i X_i'} = \cancel{\sum_{i=1}^n \l_i AX_i} + B\\
    &\iff \sum_{i=1}^n \l_i'X_i = B \iff M(t)\begin{pmatrix}\l'_1(t) \\ \vdots \\ \l'_n(t)\end{pmatrix} = B(t) \\
    &\iff \begin{pmatrix}\l_1'(t) \\ \vdots \\ \l_n'(t)\end{pmatrix} = M(t)^{-1}B(t).
\end{align*}

\section{Équations différentielles scalaires.}

\subsection{Généralités.}

\begin{defi}{}{}
    Une équation différentielle scalaire linéaire d'ordre $n$ est une équation de la forme :
    \begin{equation*}
        (E) ~:~ \forall t \in I, ~ y^{(n)} + a_{n-1}(t)y^{(n-1)} + \dots + a_0(t)y = b(t)
    \end{equation*}
    avec $a_0,...,a_{n-1},b$ des fonctions de $I$ dans $\K$.\\
    Elle est dite \bf{normalisée} car le coefficient devant $y^{(n)}$ est 1.\\Les théorèmes suivants ne s'appliquent qu'aux équations normalisées.
\end{defi}

\begin{thm}{}{}
    Soit $(E) ~:~ y^{(n)} + a_{n-1}(t)y^{(n-1)} + ... + a_0(t)y = b(t)$ pour tout $t\in I$.\\
    On suppose $a_{0},...,a_{n-1},b$ continues de $I$ dans $\K$.
    \begin{enumerate}
        \item L'ensemble des solutions de $(E_H)$ est un $\K$-espace vectoriel de dimension $n$.
        \item L'ensemble des solutions de $(E)$ est un espace affine de dimension $n$, de direction $(E_H)$.
        \item Il existe une unique solution $y$ de $(E)$ vérifiant une condition initiale donnée, de la forme :
        \begin{equation*}
            y(t_0) = \a_0, ~ y'(t_0) = \a_1, ~ \dots, ~ y^{(n-1)}(t_0) = \a_{n-1}; ~ (t_0 \in I)
        \end{equation*}
    \end{enumerate} 
    \tcblower
    \boxed{1.} Non-vide : $\zero\in S_H$, stable par combinaisons linéaires.
    \begin{equation*}
        Y=\begin{pmatrix}y\\y'\\\vdots\\ y^{n-1}\end{pmatrix}; ~ Y' = \begin{pmatrix}y'\\y''\\\vdots\\y^{(n)}\end{pmatrix}; ~ y \nt{ solution ssi } ~ Y'=AY ~ \nt{ssi} ~  A=\begin{pmatrix}0&1&0&\dots&\dots&0 \\ 0&0&1&0&\dots&0 \\ \vdots & \vdots & \ddots & \ddots & & \vdots \\ 0 & \dots & \dots & \dots & 0 & 1 \\ -a_0 & -a_1 & \dots & \dots & \dots & -a_{n-1} \end{pmatrix}
    \end{equation*}
    Notons $\tilde{S}$ l'ensemble des solutions de $Y'=AY$, on sait que $\tilde{S}$ est un $\K$-ev de dimension $n$.\\
    Posons $\phi:y\mapsto Y$, de $S_H$ dans $\tilde{S}$ Alors $\phi$ bien définie, linéaire et bijective, donc $\dim S_H = \dim\tilde{S}=n$.\\
    \boxed{3.} Le théorème de Cauchy s'applique sur $Y'=AY$, donc $\exists!Y\in\tilde{S}, ~ Y(t_0) = (\a_i)_i\in\M_{n,1}(\K)$ puis $\exists!y\in S_H$ vérifiant les conditions initiales. On a donc vérifié le théorème de Cauchy sur $(E_H)$. Pour $(E)$ :
    On a $Y'=A'(t)Y+B(t)$, donc le théorème de Cauchy s'applique, avec $A,B$ continues sur $I$.
\end{thm}

\bf{Cas particulier des équations différentielles scalaires à coefficients constants.}

On a un système différentiel autonome, on rappelle que $\X_A = \X_{A^\top} = X^n  + a_{n-1}X^{n-1} + ... + a_1X+a_0$.\\
On appelle équation caractéristique de $(E_H)$ l'équation $r^n + a_{n-1}r^{n-1} + ... + a_1r + a_0 = 0$.\\
Puis $r$ est solution de l'équation caractéristique ssi $\X_A(r)=0$ ssi $r$ est vp de $A$.\\
Si $\X_A$ est scindé à racines simples : on note $r_1,...,r_n$ ses vp distinctes, on sait qu'un S.F. de solutions de $Y'=AY$ est donné par $(X_1,...,X_n)$ où $\forall i \in \lb1,n\rb, ~ X_i = e^{r_it}V_i$ avec $(V_1,...,V_n)$ une base de $\vp$ de $A$ associés à $(r_1,...,r_n)$.\\
Puis $Y$ est solution de $Y'=AY$ ssi $\exists(\a_1,...,\a_n)\in\K^n$ tels que $Y=\a_1X_1+...+\a_nX_n$.\\
Donc $Y=\a_1e^{r_1t}V_1 + ... + \a_ne^{r_nt}V_n$ puis $y=\a_1e^{r_1t}v_1+...+\a_ne^{r_nt}v_n$ en filtrant selon la première ligne.\\
Finalement, $y$ est solution de $(E_H)$ implique $y$ est combinaison linéaire de $(e^{r_1t},...,e^{r_nt})$ donc c'est une famille génératrice de $S_H$, c'est donc une base par cardinalité.

\subsection{Le cas \texorpdfstring{$n=2$}{Lg}.}

On considère une équation différentielle scalaire linéaire d'ordre 2 : $(E) ~ : ~ y'' + a(t)y' + b(t)y = c(t)$ avec $a,b,c\in\cursive{C}(I,\K)$.\\
On appelle solution de $(E)$ une fonction $y$ deux fois dérivable sur $I$ telle que $\forall t \in I, ~ y''(t)+a(t)y'(t)+b(t)y(t)=c(t)$.\\
Pour trouver un S.F. de solutions, il suffit de trouver deux solutions de $(E_H)$ qui forment une famille libre.\\
Le théorème de Cauchy s'applique (car $a,b,c$ continues sur $I$) : il existe une unique solution de $(E)$ vérifiant une condition initiale $y(t_0)=\a$ et $y'(t_0)=\b$ avec $t_0\in I$ et $\a,\b\in\K$.\\
L'ensemble des solutions de $(E)$ est un espace affine de dimension 2, de direction $S_H$.\\
Si $(y_1,y_2)$ est une famille de solutions de $(E_H)$, on définit le Wronskien par $W(t)=\det(Y_1(t),Y_2(t))=y_1(t)y_2'(t)-y_1'(t)y_2(t)$.

\bf{Écriture matricielle.}

Pour $Y=\begin{pmatrix}y\\y'\end{pmatrix}$, $Y'=\begin{pmatrix}y'\\y''\end{pmatrix}$, on a $Y$ solution de $(E)$ ssi $\begin{pmatrix}y'\\y''\end{pmatrix}=\begin{pmatrix}0&1\\-b(t)&-a(t)\end{pmatrix}\begin{pmatrix}y\\y'\end{pmatrix}+\begin{pmatrix}0\\c(t)\end{pmatrix}$.

\begin{prop}{}{}
    Soit $(E) ~:~ y''+a(t)y'+b(t)y = c(t)$ avec $a,b,c$ continues de $I$ dans $\K$ et $(E_H) ~:~ y''+a(t)y'+b(t)y=0$.\\
    Soient $(y_1,y_2)$ une famille de solutions de $(E_H)$. Les trois propositions sont équivalents :
    \begin{enumerate}
        \item $(y_1,y_2)$ est un système fondamental de solutions de $(E)$.
        \item $\forall t \in I, ~ W(t) \neq 0$.
        \item $\exists t_0 \in I, ~ W(t_0) \neq 0$.
    \end{enumerate}
\end{prop}

\bf{Cas particulier où $a$ et $b$ sont constants.}\n
Soit $(E) ~:~ y''+ ay' + by=c(t)$ avec $c$ continue sur $I$ et $a,b\in\K$. On considère : $r^2+ar+b=0$.\\
$\bullet$ Lorsque $\K=\R$ :\\
\hspace*{0.2cm} $\hra$ Si l'équation admet deux racines réelles $r_1,r_2$, alors $(e^{r_1t}, e^{r_2t})$ est un S.F. de solutions.\\
\hspace*{0.2cm} $\hra$ Si l'équation admet une racine double réelle $r$, alors $(e^{rt},te^{rt})$ est un S.F. de solutions.\\
\hspace*{0.2cm} $\hra$ Si elle admet deux racines complexes conjuguées $\a\pm i\b$, alors $(e^{\a t}\cos(\b t), e^{\a t}\sin(\b t))$ est un S.F. de solutions.\\
On retient les deux premiers cas uniquement lorsque $\K=\C$.

Pour résoudre l'équation complète, on pourra appliquer la méthode de variation des constantes.

\begin{prop}{Variation des constantes.}{}
    Soit $(E) ~:~ y''+a(t)y'+b(t)y=c(t)$ avec $a,b,c$ continues de $I$ dans $\K$.\\
    Notons $y_1,y_2$ un système fondamental de solutions. Posons $y(t)=\l_1(t)y_1(t)+\l_2(t)y_2(t)$ pour $t\in I$.\\
    Alors $y$ est dérivable sur $I$ ssi $\l_1$ et $\l_2$ le sont, puis :
    \begin{equation*}
        y \nt{ est solution de } (E) \iff \begin{cases}\l_1' y_1 + \l_2' y_2 = 0 \\ \l_1' y_1' + \l_2'y_2' = c\end{cases}
    \end{equation*}
    Le système précédent est un système de Cramer qui détermine $\l_1'$ et $\l_2'$ de manière unique.\\
    On calcule ensuite $\l_1$ et $\l_2$ à une constante près avec deux primitivations.
\end{prop}

\begin{exercice}{}{}
    Résoudre $y'' - 5y' + 6y = te^t$ sur $\R$.
    \tcblower
    C'est une équation différentielle scalaire d'ordre 2 à coefficients continus sur $\R$. On sait que l'espace des solutions est un espace affine de dimension 2. On résout l'équation homogène : $y''-5y'+6y = 0$.\\
    É.C. : $r^2-5r+6=0$ donc $(r-2)(r-3)=0$, donc un S.F. de solutions est donné par $(e^{2t}, e^{3t})$.\\
    Solution particulière : de la forme $y=(at+b)e^{t}$.
    \begin{equation*}
        y' = (at+a+b)e^t \et y'' = (at+ 2a +b)e^{t} \ra (2at - 3a + 2b)e^t= te^t
    \end{equation*}
    Donc $a=1/2$ et $b = 3/4$. Donc $y$ solution ssi $\exists (\l,\mu)\in\R^2, ~ y=\l e^{2t} + \mu e^{3t} + \left( \frac{t}{2} + \frac{3}{4} \right)e^t$.\\
    \bf{Ou par variation des constantes :} on pose $y=\l e^{2t} + \mu e^{3t}$ pour tout $t\in\R$, dérivable ssi $\l$ et $\mu$ le sont.
    \begin{equation*}
        \begin{cases}
            \l' e^{2t} + \mu' e^{3t} = 0 \\ \l' 2e^{2t} + \mu' 3e^{3t} = te^t
        \end{cases} \iff \begin{pmatrix}\l' \\ \mu'\end{pmatrix} = \begin{pmatrix}e^{2t} & e^{3t} \\ 2e^{2t} & 3e^{3t}\end{pmatrix}^{-1}\begin{pmatrix}0\\te^t\end{pmatrix} = \begin{pmatrix}-te^{-t} \\ te^{-2t}\end{pmatrix}
    \end{equation*}
    Donc $\l'(t) = -te^{-t}$ et $\mu'(t) = te^{-2t}$ puis $\l(t) = (t+1)e^{-t} + \a$ et $\mu(t)=-t=-(\frac{t}{2} - \frac{1}{4})e^{-2t} + \b$.\\
    On en déduit $y=\l e^{2t} + \mu e^{3t} = \a e^{2t} + \b e^{3t} + (\frac{t}{2} + \frac{3}{4})e^t$.
\end{exercice}

\begin{exercice}{}{}
    Résoudre $y''-2y'+y=(t+2)e^t$ sur $\R$.
    \tcblower
    C'est une équation différentielle scalaire d'ordre 2 à coefficients continus sur $\R$. On sait que l'espace des solutions est un espace affine de dimension 2. On résout l'équation homogène : $y''-2y'+y = 0$.\\
    É.C. : $r^2-2r+1=0$ donc $(r-1)^2=0$, donc un S.F. de solutions est donné par $(e^t, te^t)$.\\
    Puisque 1 est racine double, on a $y(t)=t^2(at+b)e^t$ une solution particulière.
    \begin{equation*}
        y' = (at^3 + 3at^2 + bt^2 + 2bt)e^t \et y'' = (at^3 + 6at^2 + bt^2 + 6at + 4bt + 2b)e^t \ra (6at+2b)e^t = (t+2)e^t
    \end{equation*}
    Donc $6a=1$ et $2b=2$ donc $a=\frac{1}{6}$ et $b=1$. Donc $y$ solution ssi $\exists (\l,\mu)\in\R^2, ~ y=\l e^{t} + \mu te^{t} + \left( \frac{t^3}{6} + t^2 \right)e^t$.
\end{exercice}

\begin{exercice}{}{}
    Résoudre $y'' - 2y' + y = \frac{e^t}{1+t^2}$ sur $\R$.
    \tcblower
    C'est une équation différentielle scalaire d'ordre 2 à coefficients continus sur $\R$. On sait que l'espace des solutions est un espace affine de dimension 2. On ne cherche qu'une solution particulière.\\
    Par variation des cosntantes, on pose $y=\l e^t + \mu te^t$, dérivable ssi $\l$ et $\mu$ le sont sur $\R$.
    \begin{equation*}
        \begin{cases}
            \l' e^t + \mu' t e^t = 0 \\ \l' e^t + \mu'(t+1)e^t = \frac{e^t}{1+t^2} 
        \end{cases} \iff \begin{pmatrix}\l'\\\mu'\end{pmatrix} = \begin{pmatrix}e^t & te^t \\ e^t & (t+1)e^t\end{pmatrix}^{-1}\begin{pmatrix}0\\\frac{e^t}{1+t^2}\end{pmatrix} = \begin{pmatrix}-\frac{t}{1+t^2} \\ \frac{1}{1+t^2}\end{pmatrix}
    \end{equation*}
    Donc $\l = -\frac{1}{2}\ln(1+t^2)+\a$ et $\mu = \arctan(t) + \b$ puis $y=\b e^t + \a te^t + (t\arctan t - \frac{1}{2}\ln(1+t^2))e^t$.
\end{exercice}

\begin{exercice}{}{}
    Résoudre $x^2y'' + 3xy' + y = 0$ sur $\R_+^*$ en posant $x=e^t$.
    \tcblower
    C'est une équation différentielle scalaire linéaire d'ordre 2, homogène à coefficients continus sur $\R_+^*$. On sait que l'ensemble des solutions forme un $\R$-espace vectoriel de dimension 2 et que le théorème de Cauchy s'applique.\\
    Soit $x>0$, puisque $\exp$ est bijection de $\R$ sur $\R_+^*$, il existe un unique $t\in\R$ tel que $x=\exp(t)$.\\
    Ainsi, $y(x)=y(e^t)=z(t)$, on a $z'(t)=e^ty'(e^t)$ et $z''(t)=e^ty'(e^t)+e^{2t}y''(e^t)$.\\
    Puis $y$ est solution donc $\forall t \in \R, ~ e^{2t}y''(e^t)+3e^ty'(e^t)+y(e^t)=0$ puis $z''(t)+2z'(t)+z(t)=0$.\\
    Le changement de variable a permis d'obtenir une équaiton différentielle d'ordre 2 à coefficients constants.\\
    É.C. : $r^2+2r+1=0$ donc $(r+1)^2=0$, donc un S.F. de solutions est donné par $(e^{-t}, te^{-t})$.\\
    Puis $y(x)=y(e^t)=z(t)=z(\ln x)=\l e^{-\ln x} + \mu \ln x e^{-\ln x}$ donc $y(x)=\frac{\l}{x} + \frac{\mu \ln x}{x}$; $y_1(x)=\frac{1}{x}$ et $y_2(x)=\frac{\ln x}{x}$.\\
    On vérifie que le Wronskien ne s'annule pas : $W(x)=y_1(x)y_2'(x)-y_2(x)y_1'(x)=\frac{1}{x^3}\neq0$ sur $\R_+^*$.
\end{exercice}

\pagebreak

\begin{exercice}{}{}
    Soit $(E) ~:~ (t^2+2t+2)y'' - 2(t+1)y' + 2y = 0$.
    \begin{enumerate}
        \item Chercher des solutions polynomiales de $E$.
        \item Proposer un système fondamental de solutions et calculer le Wronskien.
    \end{enumerate}
    \tcblower
    C'est une équation différentielle scalaire linéaire d'ordre 2, homogène à coefficients continus et normalisée.\\
    Cherchons une solution polynomiale $y=\sum_{k=0}^Na_kX^k$ avec $a_N\neq0$ et $N\in\N$.\\
    On a $(t^2+2t+2)y'' - 2(t+1)y' + 2y$ est un polynôme de degré inférieur à $N$.\\
    Coefficient dominant : $N(N-1)a_N - 2Na_n + 2a_N = 0$ puis $(N-1)(N-2)=0$ : $N\in\{1,2\}$.\\
    On a donc une \bf{contrainte sur $N$}, on peut poser $y=at^2+bt+c$. Donc $y$ solution ssi $2a-b+c = 0$.\\
    Par exemple, $y_1 = t^2 - 2$ et $y_2=t+1$ forme un système fondamental de solutions de $(E)$.\\
    Puis $W(t)=y_1(t)y_2'(t)-y_2(t)y_1'(t)=-(t^2+2t+2)\neq0$.
\end{exercice}

\begin{exercice}{}{}
    Soit $(E) ~:~ ty'' + 2y' - ty = 0$.
    \begin{enumerate}
        \item Trouver une solution $y_1$ développable en série entière de $(E)$.
        \item En posant $y=y_1z$ pour $t>0$, montrer que $\sh(t)z''(t)+2\ch(t)z'(t)=0$. (Méthode de Lagrange)
        \item Déterminer $z$ et résoudre $(E)$ sur $\R_+^*$.
    \end{enumerate}
    \tcblower
    C'est une équation différentielle scalaire linéaire d'ordre 2, homogène à coefficients continus et non normalisée.\\
    L'équation est normalisée sur $]-\infty,0[$ et sur $]0,+\infty[$, donc le théorème de Cauchy s'y applique sur ces intervalles.\\
    Cherchons une solution développable en série entière. On pose $y(t)=\sum_{n=0}^{\infty}a_nt^n$ de rayon de convergence $R\geq0$.\\
    On a $y$ de classe $\cursive{C}^{\infty}$ sur $]-R,R[$, on dérive terme-à-terme avec conservation du rayon :
    \begin{equation*}
        y(t) = \sum_{n=0}^{\infty}a_nt^n; \quad y'(t) = \sum_{n=1}^\infty na_nt^{n-1}; \quad y''(t) = \sum_{n=2}^\infty n(n-1)a_nt^{n-2}
    \end{equation*} 
    Donc :
    \begin{equation*}
        \sum_{n=1}^{\infty}(n+1)na_{n+1}t^n + 2\sum_{n=0}^{\infty}(n+1)a_{n+1}t^n - \sum_{n=1}^{\infty}a_{n-1}t^n = 0
    \end{equation*}
    Par unicité du DSE : $2a_1 = 0$, $\forall n \geq 1, ~ (n+2)(n+1)a_{n+1} = a_{n-1}$.\\
    Donc $\forall n \in \N, ~ a_{2n+1} = 0$ et $a_{2n} = \frac{a_0}{(2n+1)!}$ (récurrence immédiate) donc $y=a_0\sum_{n=0}^{\infty}\frac{t^{2n}}{(2n+1)!}=a_0\frac{\sh(t)}{t}$ (RCV $\infty$).\\
    Recollement en zéro avec $y(0)=a_0$; $y$ devient $\cursive{C}^\infty$ sur $\R$ tout entier.
\end{exercice}

\pagebreak

\section{Étude qualitative.}

On cherche à obtenir des informations sur les équations différentielles sans forcément les résoudre.

\begin{exercice}{}{}
    \begin{enumerate}
        \item Résoudre $y'' + y = f(t)$ pour $t\in\R$, où $f:\R\to\R$ continue.
        \item Soit $f:\R\to\R$ de classe $\cursive{C}^2$ sur $\R$ telle que $f(x)+f''(x)\geq0$ pour tout $x\in\R$.\\ Montrer que $f(x) + f(x+\pi)\geq0$ pour tout $x\in\R$.
    \end{enumerate}
    \tcblower
    \boxed{1.} On reconnaît une équation différentielle scalaire d'ordre 2 linéaire à coefficients continus. On sait que l'ensemble des solutions forment un espace affine de dimension 2 et que le théorème de Cauchy s'applique :
    On sait que $y$ est solution de l'équation homogène ssi $\exists (\l,\mu)\in\R^2, ~ \forall t \in \R, ~  y =\l\cos t + \mu \sin t$.\\
    Avec variation des constantes : on pose $y(t) = \l(t)\cos(t) + \mu(t)\sin(t)$, $y$ est dérivable ssi $\l$ et $\mu$ le sont.\\
    Puis $y$ est solution ssi :
    \begin{equation*}
        \begin{pmatrix} \cos t & \sin t \\ -\sin t & \cos t \end{pmatrix}  \begin{pmatrix}\l ' \\ \mu'\end{pmatrix} = \begin{pmatrix}0 \\ f(t)\end{pmatrix} \iff \begin{pmatrix}\l' \\ \mu '\end{pmatrix} = \begin{pmatrix}\cos t & -\sin t \\ \sin t & \cos t\end{pmatrix}\begin{pmatrix}0 \\ f(t) \end{pmatrix}
    \end{equation*}
    Donc $\l'= -\sin(t)f(t)$ et $\mu' = \cos(t)f(t)$, puis $\l(t) = \int_0^t - \sin(u)f(u)\du+\a$ et $\mu(t)=\int_0^t\cos(u)f(u)\du+\b$.
    \begin{align*}
        y &= \l(t)\cos(t) + \mu(t)\sin(t) = \a\cos t + \b \sin t + \int_0^t (-\cos t \sin u + \cos u \sin t)f(u)\du \\&= \a\cos t + \b \sin t + \int_0^t \sin(t-u)f(u)\du.
    \end{align*}
    \boxed{2.} Soit $f\in\cursive{C}^2(\R)$ telle que $f(x)+f''(x)\geq0$ pour tout $x\in\R$.\\
    On pose $\phi(x) = f(x) + f''(x)\geq0$ pour $x\in\R$, alors $\phi$ est continue sur $\R$ et positive, $f$ est solution de $y''+y=\phi$.\\
    D'après le 1, on a :
    \begin{equation*}
        \exists(\a,\b)\in\R^2, ~ \forall t \in \R, ~ f(t)=\a\cos t + \b\sin t + \int_0^t \sin(t-u)\phi(u)\du
    \end{equation*}
    Soit $x\in\R$. \small
    \begin{align*}
        f(x)+f(x+\pi)&=\a(\cos x + \cos(x + \pi)) + \b(\sin x + \sin(x+\pi)) + \int_0^x\sin(x-t)\phi(t)\dt + \int_0^{x+\pi}\sin(x+\pi-t)\phi(t)\dt\\
        &=\int_0^x \sin(x-t)\phi(t)\dt - \int_0^{x+\pi} \sin(x-t)\phi(t)\dt = \int_x^{x+\pi} \sin(t-x)\phi(t)\dt
    \end{align*}
    Avec le changement de variable $u=t-x$ : $f(x)+f(x+\pi)=\int_0^\pi \sin u \phi(u+x)\du \geq 0$ comme intégrale de fonction positive.
\end{exercice}

\pagebreak

\begin{exercice}{}{}
    Soit $(E) ~:~ y'' + a(t)y'(t) + b(t)y = 0$ avec $a,b$ continues sur $I$.
    \begin{enumerate}
        \item Soit $t_0 \in I$, montrer que $y(t_0) = y'(t_0) = 0 \ra y=0$ sur $I$.
        \item Soient $t_0<t_1 \in I$, montrer que $y=0$ sur $[t_0,t_1] \ra y=0$ sur $I$.
        \item Soient $t_0<t_1\in I$, montrer que $y$ s'annule un nombre fini de fois sur $[t_0,t_1]$ si $y\neq0$.
        \item Montrer que les zéros de $y\neq0$ sont isolés. Justifier que l'on peut parler de << zéros consécutifs >> de $y$.
    \end{enumerate}
    \tcblower
    On reconnaît une équation différentielle scalaire linéaire d'ordre 2 homogène à coefficients continus, normalisée. L'ensemble des solutions de $(E)$ est un $\R$-ev de dimension 2 et que le théorème de Cauchy s'applique.\\
    \boxed{1.} Cauchy : $\forall (\a,\b)\in\R^2, ~ \forall t_0 \in I$, il existe une unique solution $y$ de $E$ tq $y(t_0)=\a$ et $y'(t_0)=\b$.\\
    Ici, $\a=\b=0$, donc il existe une unique solution $y$ de $E$ telle que $y(t_0)=y'(t_0)=0$, la fonction $\0$ convient.\\
    \boxed{2.} Soit $t_2 \in ]t_0,t_1[$. Puisque $y=0$ sur $[t_0,t_1]$, $y(t_2)=y'(t_2)=0$. D'après $1$, $y=0$ sur $I$.\\
    \boxed{3.} Par l'absurde, supposons que $y$ s'annule une infinité de fois sur $[t_0,t_1]$.\\
    Notons $(\a_n)_n$ une suite de $[t_0,t_1]$ d'éléments tous distincts telle que $\forall n \in \N, ~ y(\a_n)=0$.\\
    On peut extraire de $(\a_n)_n$ une suite convergente par B-W : $\exists \phi : \N \to \N$ telle que $\a_{\phi(n)}\to \a \in [t_0,t_1]$.\\
    Par continuité de $y$ sur $I$ donc en $\a$ : $\a_{\phi(n)}\to \a$ implique $y(\a)=0$ car $\forall n \in \N, ~ y(\a_{\phi(n)})=0$.\\
    On applique le théorème de Rolle sur chacun des intervalles $(\a_{\phi(n)}, \a_{\phi(n+1)})$.\\
    Ainsi, $\exists \b_n \in ]\a_{\phi(n)}, \a_{\phi(n+1)} [$ tel que $y'(\b_n)=0$ par Rolle. (on a supposé spdg $\a_{\phi(n)} < \a_{\phi(n+1)}$).\\
    Puis par encadrement, $\b_n \to \a$, donc par continuité de $y'$ en $\a$ : $y'(\a)=0$ puisque $\forall n \in \N, ~ y'(\b_n)=0$.\\
    Donc $y(\a)=0$ et $y'(\a)=0$, donc $y=0$ par 1, absurde.\\
    \boxed{4.} Soit $\a$ un zéro de $y$, montrons que $\exists h > 0, ~ \forall t \in ]\a - h, \a + h[ \setminus\{\a\}, ~ y(t)\neq0$.\\
    Puisque $y$ n'est pas nulle, $y'(\a)\neq0$ par 1, SPDG, $y'(\a)>0$.\\
    Puis $y'$ est continue sur $I$ donc $\exists h > 0, ~ \forall t \in ]\a-h,\a+h[, ~  y'(t)>0$.\\
    Puis $y$ est strictement croissante sur $]\a-h,\a+h[$ et $\a$ est le seul zéro de $y$ sur $]\a-h,\a+h[$.\\
    Donnons un sens à la notion de << zéros consécutifs >> :\\
    Soit $\a$ un zéro de $y$. On suppose que $y$ possède un zéro $\b>\a$, $\a$ est isolé par ce qui précède : $\exists h > 0$ tel que $y$ ne s'annule pas sur $]\a,\a+h[$.\\
    On pose $E=\{t>a+h, ~ y(t)=0\}$, non-vide (contient $\b$) et minoré par $\a+h$, donc possède une borne inf $\tilde{a}$.\\
    Alors $\tilde{a}$ est la limite d'une suite $(t_n)_n$ de $E$; par continuité de $y$ en $\tilde{a}$ : $y(t_n)\to y(\tilde{\a})=0$.\\
    Par construction, $\tilde{a}$ est le plus petit zéro de $y$ après $\a$, alors $\a$ et $\tilde{\a}$ sont dits zéros consécutifs.
\end{exercice}

\begin{lemme}{}{}
    Une fonction convexe ou concave sur $\R$, de classe $\cursive{C}^1$ et bornée est constante.
    \tcblower
    Soit $f$ convexe sur $\R$, de classe $\cursive{C}^1$ et bornée. Supposons qu'il existe $t_0\in\R, ~ f'(t_0)\neq0$.\\
    $\hra$ Si $f'(t_0)>0$, alors $\forall t \in \R, ~ f(t) \geq f(t_0) + (t-t_0)f'(t_0)$ donc $f(t)\xrightarrow[t\to+\infty]{}+\infty$, absurde.\\
    $\hra$ Si $f'(t_0)<0$, alors $\forall t \in \R, ~ f(t) \geq f(t_0) + (t-t_0)f'(t_0) \xrightarrow[t\to-\infty]{}+\infty$, absurde. 
\end{lemme}

\pagebreak

\begin{exercice}{Équation de Sturm-Liouville.}{}
    Soit $(E) ~:~ y'' + a(t)y = 0$ pour tout $t\in I$, avec $q$ continue sur $I$.
    \begin{enumerate}
        \item Soit $y_1,y_2$ un S.F. de $E$ de Wronskien $W$, montrer que $W$ est constant sur $I$.
        \item On suppose $q\geq0$ sur $\R$ et non identiquement nulle. Soit $y$ une solution de $(E)$ bornée.\\
        Montrer que $y$ s'annule au moins une fois.
        \item On prend $q(t)=-e^t$ sur $\R$ et $y$ une solution de $(E)$. Étudier la convexité de $y^2$.\\
        En déduire $y(0)=y(1)=0 \ra y = 0$.
        \item Soit $y_1,y_2$ un S.F. de $E$, et $a,b$ deux zéros consécutifs de $y_1$. Montrer que $y_2$ s'annule exactement une fois sur $]a,b[$.
    \end{enumerate}
    \tcblower
    Une équation de Sturm-Liouville est une équation différentielle liénaire d'ordre 2 normalisée scalaire homogène à coefficients continus. L'ensemble de solutions est un $\R$-ev de dimension 2. Ce théorème de Cauchy s'applique.\\
    \boxed{1.} $W=y_1y'_2-y'_1y_2$ est dérivable sur $I$ et $W'=y_1y''_2-y''_1y_2$.\\
    Par hypothèse, $y''_1+qy_1 = 0$ et $y''_2 + qy_2 = 0$ donc $W' = y_1(-qy_2)-(-qy_1)y_2=0$.\\ Donc $W$ est constant sur l'intervalle $I$.\\
    \boxed{2.} Par hypothèse, $y''+qy=0$, supposons par l'absurde que $y$ ne s'annule pas sur $I$.\\
    Supposons $y$ continue sur $\R$ et ne s'annule pas, donc par TVI, $y$ garde un signe constant sur $\R$.\\
    SPDG, $y>0$ sur $\R$ puis $y'' = -qy \leq 0$ sur $\R$ donc $y'$ est concave sur $\R$.\\
    Puisque $y$ est concave sur $\R$, de classe $\cursive{C}^2$ et bornée, elle est constante, donc $y''=0$ sur $\R$.\\
    Or $y''=-qy$ donc $qy=0$ donc $q=0$ puisque $y>0$, absurde.\\
    \boxed{3.} Posons $z=y^2$, deux fois dérivables : $z' = 2yy'$ et $z''=2y'^2 + 2yy' = 2(y')^2 + 2e^xy\geq0$.\\
    Donc $z$ est convexe sur $\R$ donc $\forall t \in [0,1], ~ z(t)\leq 0$ (corde entre 0 et 1). Or $z\geq0$ (comme carré).\\
    Donc $\forall t \in [0,1], ~ z(t)=0$.\\
    \boxed{4.} Posons $W=y_1y'_2-y'_1y_2$, on a déjà vu que $W$ est constant ($W'=0$) et $W(a)=-y_1'(a)y_2(a)$.\\
    SPDG, $y_1\geq0$ sur $[a,b]$, alors $y'_1(a)\geq0$ car $y_1'(a)=\lim \frac{y_1(t)-y_1(a)}{t-a}\geq0$.\\
    De plus, $y'_1(b)\leq0$. Puis $y_1'(a)>0$ (car sinon $y_1(a)=y'_1(a)=0$ puis $y_1=0$ absurde) et $y_1'(b)<0$.\\
    Donc $W(a)=-y_1'(a)y_2(a)$ et $W(b)=-y_1'(b)y_2(b)$ et $W(a)=W(b)$ car $W$ est constant.\\
    Donc $y_2(a)$ et $y_2(b)$ sont de signes contraires puisque $y'_1(a)$ et $y'_1(b)$ le sont (non nuls car $W\neq0$).\\
    On a $y_2$ continue sur $[a,b]$ donc $y_2$ s'annule au moins une fois sur $[a,b]$ par TVI.\\
    Supposons par l'absurde que $y_2$ s'annule deux fois sur $]a,b[$ : en $c$ et en $d$, alors en échangeant les rôles de $y_1$ et $y_2$, $y_1$ s'annule au moins une fois entre $c$ et $d$, c'est absurde car $a$ et $b$ sont zéros consécutifs de $y_1$.
\end{exercice}

\begin{exercice}{}{}
    Soit $(E)$ l'ensemble des applications dérivables $f$ sur $\R$ telles que $f'(t) = f(-t)$ pour tout $t\in\R$.
    \begin{enumerate}
        \item Vérifier que $E$ est un $\R$-ev.
        \item Résoudre et donner une base de $(E)$.
    \end{enumerate}
    \tcblower
    \boxed{1.} C'est un sous-espace vectoriel de $\cursive{C}^1(\R,\R)$.\\
    $\hra$ La fonction nulle est solution.\\
    $\hra$ Soient $f,g\in E$ et $\l\in\R$, alors $(f+\l g)'(t) = f'(t) + \l g'(t) = f(-t) + \l g(-t) = (f + \l g)(-t)$.\\
    Donc $f+\l g \in E$.\\
    \boxed{2.} Déjà, $f\in E \ra f\in\cursive{C}^2(\R)$ car $f'=f\circ (-\Id)$, et même $\cursive{C}^\infty(\R)$.\\
    Soit $f$ solution. On dérive l'équation : $f''(t) = -f'(-t) = -f(t)$ donc $f''(t) + f(t) = 0$.\\
    Donc $\exists \l,\mu\in\R, ~ f(t)=\l\cos t + \mu\sin t$.\\
    \bf{Synthèse.} Soit $f(t) = \l \cos t + \mu \sin t$ pour tout $t\in\R$, avec $(\l,\mu)\in\R^2$.\\
    Alors $f$ est dérivable, et $f'(t)=-\l \sin t + \mu \cos t$ puis $\forall t \in \R, ~ f'(t)=f(-t) \iff \l = \mu$ avec $(t=0)$.\\
    Donc $E=\{t\mapsto \l(\cos t + \sin t), ~ \l \in \R\}=\Vect(t\mapsto \cos t + \sin t)$ et $\dim E = 1$.
\end{exercice}

\begin{exercice}{}{}
    Soit $(E)~:~X'=AX+B$ avec $A:I\to\M_n(\R)$ et $B:I\to\M_{n,1}(\R)$ continues sur $I$ intervalle de $\R$.\\
    On note $W$ le Wronskien associé à un SF de solutions $(X_1,...,X_n)$.\\
    Montrer que $W$ est solution de l'équation différentielle $W'=\Tr(A)W$.
    \tcblower
    Par définition, $W(t)=\det(X_1(t),...,X_n(t))$ pour tout $t\in I$ et $(X_1,...,X_n)$ est base de $S_H$.\\
    Donc pour tout $i\in\lb1,n\rb, ~ X_i$ est dérivable sur $I$ et $X'_i = AX_i$.\\
    Puis $\det$ est $n$-linéaire et $X_1,...,X_n$ dérivables sur $I$ donc $W$ est dérivable sur $I$ et :
    \begin{align*}
        W'(t) &= \det(X'_1, X_2, \dots, X_n) + \det(X_1, X'_2, \dots, X_n) + \dots + \det(X_1,X_2, \dots, X'_n)\\
        &= \det(AX_1, X_2, \dots, X_n) + \det(X_1, AX_2, \dots, X_n) + \dots + \det(X_1,X_2,\dots,AX_n).
    \end{align*} 
    Posons $\phi:(Y_1,...,Y_n) \mapsto \det(AY_1,...,Y_n) + ... + \det(Y_1,...,AY_n)$.\\
    L'ensemble des formes $n$-linéaires alternées est un $\R$-espace vectoriel de dimension 1.\\
    Vérifions que $\phi$ est une forme $n$-linéaire alternée.\\
    $\hra$ Linéaire par rapport à la première variable : soient $y_1,z_1 \in \M_{n,1}(\R)$ et $\l\in\R$.
    \begin{align*}
        \phi(\l y_1 + z_1, ..., y_n) &= \det(\l Ay_1 + Az_1, ..., y_n) + ... + \det(\l y_1 + z_1, ... + Ay_n)\\
        &=\l\det(Ay_1+z_1, ..., y_n) + \det(Az_1, ..., y_n)+ ... + \l\det(y_1,...,Ay_n) + \det(z_1,...,Ay_n) \\
        &= \l \phi(y_1,...,y_n) + \phi(z_1,...,y_n).
    \end{align*}
    Puis linéaire de la même façon selon les autres variables.\\
    $\hra$ Alternée, par exemple pour $y_1=y_2$ :
    \begin{align*}
        \phi(y_1,...,y_n) &= \det(Ay_1, y_2, ..., y_n) + det(y_1,Ay_2,...,y_n) + ... + \det(y_1,y_2,...Ay_n)\\
        &=\det(Ay_1,y_1,...,y_n) + \det(y_1,Ay_1, ..., y_n) + ... + \det(y_1,y_1,...,Ay_n)\\
        &=\det(Ay_1,y_1,...,y_n) - \det(Ay_1,y_1, ..., y_n) = 0
    \end{align*}
    Donc $\phi$ est $n$-linéaire alternée donc $\exists \l \in \R, ~ \phi = \l \det$ donc :
    \begin{equation*}
        \forall t \in I, ~ W'(t) = \phi(X_1,...,X_n) = \l\det(X_1,...,X_n) = \l W(t).
    \end{equation*}
    Soient $(E_1,...,E_n)$ les vecteurs de la base canonique de $\R^n$.\\
    Alors $\phi(E_1,...,E_n) = \det(AE_1, E_2, ..., E_n) + ... + \det(E_1,..., AE_n)$ et pour tout $i\in\lb1,n\rb$ :
    \begin{equation*}
        \det(E_1,..., AE_i, ..., E_n) = \begin{vmatrix}1 & 0 & \dots & a_{1,i} & \dots & 0 \\ 0 & 1 & \dots & a_{2,i} & \dots & 0 \\ \vdots & \vdots & & \vdots & & \vdots \\ 0 & 0 & \dots & a_{n,i} & \dots & 1\end{vmatrix} = a_{i,i}.
    \end{equation*}
    Donc $\l=\sum_{i=1}^n a_{i,i}=\Tr(A)$ puis $W'(t)=\Tr(A(t))W(t)$.
\end{exercice}

\pagebreak

\begin{exercice}{}{}
    \begin{enumerate}
        \item Soit $A\in\A_n(\R)$. Montrer que $\exp(A)\in O_n(\R)$.
        \item Soit $S\in\M_n(\R)$ telle que $\forall U \in O_n(\R), ~ \Tr(US)\leq \Tr(S)$.
        \begin{enumerate}
            \item Soit $f(t)=\Tr(\exp(tA)S)$ avec $A\in \A_n(\R)$ et $t\in\R$, montrer que $f$ est dérivable et $f'(0)=0$.
            \item En déduire que $\forall A \in \A_n(\R), ~ \Lg S,A \Rg = 0$ puis $S\in\S_n(\R)$.
            \item Montrer que $\Sp(S)\subset\R_+$.
        \end{enumerate}
        \item $\forall U \in O_n(\R), ~ \forall A \in \M_n(\R) ~:~ \phi(U)=\Tr(UA)$.
        \begin{enumerate}
            \item Montrer que $\phi$ atteint un maximum en $U_0\in O_n(\R)$.
            \item Montrer que $S=U_0A\in\S_n^+(\R)$.
            \item Retrouver la décomposition polaire.
        \end{enumerate}
    \end{enumerate}
    \tcblower
    \boxed{1.} Soit $A\in \A_n(\R)$, on a $(\exp A)^\top \exp A = \exp( A^\top)\exp(A) = \exp(-A)\exp(A)=\exp(0)=I_n:\exp(A)\in O_n(\R)$.\\
    \boxed{2.a.} On sait que $g:t\mapsto \exp(tA)$ est dérivable sur $\R$ et $g'(t)=A\exp(tA)$.\\
    Ainsi, $t\mapsto \phi(t)S$ est dérivable sur $\R$ puis $\Tr$ est linéaire donc $f$ est dérivable sur $\R$. \\
    Avec $Q1$ : $\forall t \in \R, ~ tA \in \A_n(\R)$ donc $\exp(tA)\in O_n(\R)$ donc $f(t)=\Tr(\exp(t A) S) \leq \Tr(S) = f(0)$.\\
    Donc 0 est maximum global de $f$ sur $\R$, donc $f'(0)=0$.\\
    \boxed{2.b.} On a $f'(0)=\Tr(AS)$ donc $\Tr(AS)=0$.
    \begin{equation*}
        \forall A \in \A_n(\R), ~ \Tr(AS) = 0 \iff \Tr(A^\top S) = 0 \iff \Lg A,S \Rg = 0 \iff S \in \A_n(\R)^\bot = \S_n(\R).
    \end{equation*}
    \boxed{2.c.} Théorème spectral : $\exists P \in O_n(\R), ~ \exists D \in \cursive{D}_n(\R), ~ S=PDP^\top$.
    \begin{equation*}
        \forall U \in O_n(\R), ~ \Tr(US) = \Tr(UPDP^\top) \leq \Tr(S)
    \end{equation*}
    Soit $V\in O_n(\R)$, et $U=PVP^\top\in O_n(\R)$, alors $\Tr(US)=\Tr(PVDP^\top)=\Tr(VD)$.\\
    Posons $D=\Diag(\l_1,...,\l_n)$, avec $\l_1,...,\l_p$ valeurs propres de $S$, et soit $V\in O_n(\R)$.
    \begin{equation*}
        [VD]_{i,j} = \sum_{k=1}^n v_{i,k}d_{k,j} = v_{i,j}\l_j \ra \Tr(VD) = \sum_{i=1}^n v_{i,i}\l_i.
    \end{equation*}
    Donc $\forall V \in O_n(\R), ~ \sum_{i=1}^nv_{i,i}\l_i \leq \sum_{i=1}^n \l_i$.\\
    Spécifions en $V=\Diag(\e_i)$ avec $\e_i \in \{-1,1\}$ du signe de $\l_i$.
    \begin{equation*}
        \sum_{i=1}^n v_{i,i}\l_i = \sum_{i=1}^n |\l_i| \leq \sum_{i=1}^n\l_i.
    \end{equation*}
    Donc $\forall i \in \lb1,n\rb, ~ \l_i \geq 0$.\\
    \boxed{3.a.} $O_n(\R)$ est compact de $\M_n(\R)$ puis $\phi$ y est continue et à valeurs réelles donc elle est bornée sur $O_n(\R)$ et atteint ses bornes, donc $\exists U_0 \in O_n(\R), ~ \forall U \in O_n(\R), ~ \phi(U) \leq U_0$.\\
    \boxed{3.b.} Soit $U\in O_n(\R), ~ \Tr(US)=\Tr(UU_0A) = \phi(UU_0) \leq \phi(U_0) = \Tr(S)$ donc $S\in\S_n^+(\R)$ par la partie 2.\\
    \boxed{3.c.} Soit $A\in\M_n(\R)$, on a $S=U_0A$ donc $U_0^\top S = A$ donc $A=\O S$ avec $\O = U_0^\top \in O_n(\R)$ et $S\in\S_n^+(\R)$.
\end{exercice}

\end{document}
